% GENERATED FILE. Edit canonical lessons, then run npm run generate:books.
% canonical-source-hash: fnv1a64:d1135887c795960b
% canonical-lessons: LA-C186-gubernator, LA-C186-mercatus, LA-C186-pondus, LA-C186-mensura, LA-C186-libra

\chapter{Helmsman, Market, Weight, Measure, Pound}
\label{ch:la-helmsman-market-weight-measure-pound}
\hlchaptermodality{186}

\begin{chapteropening}
\textbf{By the end of this chapter:} \emph{I can say a helmsman, a market, a weight, a measure and a pound.}

Complete the last lesson of chapter 186: I can say a helmsman, a market, a weight, a measure and a pound.
\end{chapteropening}

\section[gubernātor, gubernātōris]{gubernātor, gubernātōris — a helmsman}

\label{lesson:LA-C186-gubernator}



\begin{quote}
Before the new one: say the Latin for a rain basin, then the Latin for a foundation.
\end{quote}

\subsection*{You'll want to know: gubernātor, gubernātōris}
\textbf{gubernātor, gubernātōris} — \textquotedblleft{}a helmsman\textquotedblright{}.

\begin{cousinweb}
\textbf{gubernātor} is built on \emph{gubernō}, \textquotedblleft{}to steer\textquotedblright{} (EtymDB 2.1), a verb Latin took from Greek \emph{kybernaō}, \textquotedblleft{}to steer, to pilot\textquotedblright{}, with \emph{-tor}, \textquotedblleft{}one who does it\textquotedblright{}. English \emph{governor} comes from it through French.
\end{cousinweb}

\begin{grammarlens}[title={What you've built}]
One more word for this chapter.
\end{grammarlens}

\subsection*{Your turn}
\begin{itemize}
\raggedright
  \item \emph{Say it:} \emph{gubernātor, gubernātōris}
  \item \emph{Say it:} \emph{gubernātor, gubernātōris}, once more
  \item \emph{Say it:} say \emph{gubernātor, gubernātōris}
  \item \emph{Recall:} say the Latin for a rain basin, then the Latin for a foundation, then say \emph{gubernātor, gubernātōris} again
\end{itemize}

\begin{tcolorbox}[breakable,colback=teal!4,colframe=teal!35!black,title={Before you move on}]
What does gubernātor, gubernātōris mean? (\textquotedblleft{}A helmsman\textquotedblright{}.) Say it once more.
\end{tcolorbox}
\section[mercātus, mercātūs]{mercātus, mercātūs — a market, a fair}

\label{lesson:LA-C186-mercatus}



\begin{quote}
Before the new one: say the Latin for a foundation, then the Latin for a helmsman.
\end{quote}

\subsection*{You'll want to know: mercātus, mercātūs}
\textbf{mercātus, mercātūs} — \textquotedblleft{}a market, a fair\textquotedblright{}.

\begin{cousinweb}
\textbf{mercātus} is built on \emph{mercor}, \textquotedblleft{}to trade\textquotedblright{} (EtymDB 2.1), from \emph{merx}, \textquotedblleft{}goods\textquotedblright{}. English \emph{market} comes from it.
\end{cousinweb}

\begin{grammarlens}[title={What you've built}]
One more word for this chapter.
\end{grammarlens}

\subsection*{Your turn}
\begin{itemize}
\raggedright
  \item \emph{Say it:} \emph{mercātus, mercātūs}
  \item \emph{Say it:} \emph{mercātus, mercātūs}, once more
  \item \emph{Say it:} say \emph{mercātus, mercātūs}
  \item \emph{Recall:} say the Latin for a foundation, then the Latin for a helmsman, then say \emph{mercātus, mercātūs} again
\end{itemize}

\begin{tcolorbox}[breakable,colback=teal!4,colframe=teal!35!black,title={Before you move on}]
What does mercātus, mercātūs mean? (\textquotedblleft{}A market, a fair\textquotedblright{}.) Say it once more.
\end{tcolorbox}
\section[pondus, ponderis]{pondus, ponderis — a weight}

\label{lesson:LA-C186-pondus}



\begin{quote}
Before the new one: say the Latin for a helmsman, then the Latin for a market.
\end{quote}

\subsection*{You'll want to know: pondus, ponderis}
\textbf{pondus, ponderis} — \textquotedblleft{}a weight\textquotedblright{}.

\begin{cousinweb}
\textbf{pondus} belongs with \emph{pendō}, \textquotedblleft{}to weigh\textquotedblright{} (Lewis \& Short), and is usually traced to Proto-Indo-European \emph{*(s)pend-}, \textquotedblleft{}to stretch, to pull\textquotedblright{} (EtymDB 2.1). English \emph{pound} comes from its family.
\end{cousinweb}

\begin{grammarlens}[title={What you've built}]
One more word for this chapter.
\end{grammarlens}

\subsection*{Your turn}
\begin{itemize}
\raggedright
  \item \emph{Say it:} \emph{pondus, ponderis}
  \item \emph{Say it:} \emph{pondus, ponderis}, once more
  \item \emph{Say it:} say \emph{pondus, ponderis}
  \item \emph{Recall:} say the Latin for a helmsman, then the Latin for a market, then say \emph{pondus, ponderis} again
\end{itemize}

\begin{tcolorbox}[breakable,colback=teal!4,colframe=teal!35!black,title={Before you move on}]
What does pondus, ponderis mean? (\textquotedblleft{}A weight\textquotedblright{}.) Say it once more.
\end{tcolorbox}
\section[mēnsūra, mēnsūrae]{mēnsūra, mēnsūrae — a measure}

\label{lesson:LA-C186-mensura}



\begin{quote}
Before the new one: say the Latin for a market, then the Latin for a weight.
\end{quote}

\subsection*{You'll want to know: mēnsūra, mēnsūrae}
\textbf{mēnsūra, mēnsūrae} — \textquotedblleft{}a measure\textquotedblright{}.

\begin{cousinweb}
\textbf{mēnsūra} is built on \emph{mēnsus}, \textquotedblleft{}measured\textquotedblright{}, the participle of \emph{mētior}, \textquotedblleft{}to measure\textquotedblright{} (EtymDB 2.1), with \emph{-ūra}, an ending for an act or result.
\end{cousinweb}

\begin{grammarlens}[title={What you've built}]
One more word for this chapter.
\end{grammarlens}

\subsection*{Your turn}
\begin{itemize}
\raggedright
  \item \emph{Say it:} \emph{mēnsūra, mēnsūrae}
  \item \emph{Say it:} \emph{mēnsūra, mēnsūrae}, once more
  \item \emph{Say it:} say \emph{mēnsūra, mēnsūrae}
  \item \emph{Recall:} say the Latin for a market, then the Latin for a weight, then say \emph{mēnsūra, mēnsūrae} again
\end{itemize}

\begin{tcolorbox}[breakable,colback=teal!4,colframe=teal!35!black,title={Before you move on}]
What does mēnsūra, mēnsūrae mean? (\textquotedblleft{}A measure\textquotedblright{}.) Say it once more.
\end{tcolorbox}
\section[lībra, lībrae]{lībra, lībrae — a pound; a pair of scales}

\label{lesson:LA-C186-libra}



\begin{quote}
Before the new one: say the Latin for a weight, then the Latin for a measure.
\end{quote}

\subsection*{You'll want to know: lībra, lībrae}
\textbf{lībra, lībrae} — \textquotedblleft{}a pound; a pair of scales\textquotedblright{}.

\begin{cousinweb}
\textbf{lībra} has no sure origin; Greek \emph{litra}, \textquotedblleft{}a pound\textquotedblright{}, is often compared (Lewis \& Short), and both may come from a common Mediterranean word. The abbreviation lb for the pound comes from it.
\end{cousinweb}

\begin{grammarlens}[title={What you've built}]
One more word for this chapter.
\end{grammarlens}

\subsection*{Your turn}
\begin{itemize}
\raggedright
  \item \emph{Say it:} \emph{lībra, lībrae}
  \item \emph{Say it:} \emph{lībra, lībrae}, once more
  \item \emph{Say it:} say \emph{lībra, lībrae}
  \item \emph{Recall:} say the Latin for a weight, then the Latin for a measure, then say \emph{lībra, lībrae} again
\end{itemize}

\begin{tcolorbox}[breakable,colback=teal!4,colframe=teal!35!black,title={Before you move on}]
What does lībra, lībrae mean? (\textquotedblleft{}A pound; a pair of scales\textquotedblright{}.) Say it once more.
\end{tcolorbox}
